% Vietnamese translation for OI Public Library
% Source-File: IOI中国国家候选队论文/2005/朱泽园/朱泽园.doc
\documentclass[11pt]{article}
\usepackage{oipl-vn}

\title{Trở về điểm xuất phát: một kiểu tư duy đột phá}
\author{Zhu Zeyuan}
\date{2005}

\begin{document}
\maketitle

\origtitle{Back to the starting point: a breakthrough way of thinking}
\sourcepath{IOI China candidate team papers / 2005 / Zhu Zeyuan / Zhu Zeyuan.doc}

\section*{Tóm tắt}

Ở mức cao, thi tin học không còn chỉ là cuộc so tài giữa một thuật toán hay
một cấu trúc dữ liệu đơn lẻ, mà là phép thử đối với phương pháp tư duy và năng
lực sáng tạo. Bài viết này phân tích một kiểu tư duy đột phá: \term{trở về
điểm xuất phát}, tức là dám tạm bỏ kết quả đang có để quay lại bước phân tích
ban đầu và suy nghĩ một cách tỉnh táo.

Bài viết bắt đầu từ một bài toán thông thường vốn đã có lời giải kinh điển.
Sau khi nhắc lại lời giải đó, ta quay lại thuật toán thô ban đầu, tìm ra một
thuật toán đơn giản đã bị bỏ quên, rồi tối ưu và liên hệ nó với rừng tập rời
nhau. Cuối cùng bài viết nhấn mạnh ý nghĩa của việc tiến lên thông qua những
đoạn đường vòng trong tư duy.

\section{Đặt vấn đề}

\subsection{Mô tả bài toán}

Bài toán là một biến thể của \textit{USACO 2.1 Shaping Regions}. Có \(N\)
hình chữ nhật đục màu khác nhau, với \(1\le N\le 3000\), được đặt lên một tờ
giấy trắng có chiều dài \(A\) và chiều rộng \(B\). Mọi cạnh của các hình chữ
nhật đều song song với cạnh của tờ giấy, và mọi hình chữ nhật đều nằm trong
tờ giấy. Gốc tọa độ \((0,0)\) nằm ở góc trái dưới của tờ giấy.

\paragraph{Dữ liệu vào.}
Dòng đầu chứa \(A,B,N\). Từ dòng thứ hai đến dòng \(N+1\), các hình chữ nhật
được cho theo thứ tự từ dưới lên trên. Mỗi dòng gồm năm giá trị
\[
  \mathit{llx},\mathit{lly},\mathit{urx},\mathit{ury},\mathit{color},
\]
là tọa độ góc trái dưới, tọa độ góc phải trên và màu của hình chữ nhật. Mọi
tọa độ, kể cả \(A,B\), là số thực trong đoạn từ \(0\) đến \(10^9\). Màu \(1\)
trùng với màu trắng của nền giấy.

\paragraph{Dữ liệu ra.}
In danh sách các màu nhìn thấy được cùng diện tích tương ứng, giữ \(3\) chữ
số sau dấu thập phân. Danh sách được sắp theo thứ tự tăng của mã màu; không
in màu có diện tích bằng \(0\).

\paragraph{Ví dụ.}
\begin{verbatim}
20 20 3
2 2 18 18 2
0 8 19 19 3
8 0 10 19 4
\end{verbatim}

\begin{verbatim}
1 91.000
2 84.000
3 187.000
4 38.000
\end{verbatim}

\subsection{Phân tích ban đầu: rời rạc hóa}

Từ bài \textit{Picture} của IOI 1998, các bài toán đặt hình chữ nhật trên mặt
phẳng xuất hiện nhiều trong thi tin học. Tiền xử lý cơ bản nhất là rời rạc
hóa mặt phẳng: lấy toàn bộ tọa độ \(x\) và \(y\) từng xuất hiện trong các
hình chữ nhật, sắp xếp, loại trùng rồi đánh số theo thứ tự.

Đoạn nối hai điểm lưới kề nhau gọi là \term{đoạn nguyên tố}; vùng giữa hai
tọa độ \(x\) kề nhau gọi là \term{cột rời rạc}, vùng giữa hai tọa độ \(y\) kề
nhau gọi là \term{hàng rời rạc}; giao của một hàng và một cột rời rạc gọi là
\term{ô rời rạc}. Mỗi tọa độ đỉnh hình chữ nhật ban đầu đều tương ứng với một
tọa độ lưới.

Rời rạc hóa giúp xử lý tọa độ thực hoặc tọa độ vượt quá kiểu nguyên thông
thường. Sau rời rạc hóa, mọi tọa độ đỉnh hình chữ nhật có thể biểu diễn bằng
số nguyên từ \(1\) đến \(2N\). Nhờ bảng tra hoặc băm, việc chuyển đổi giữa
tọa độ thật và tọa độ rời rạc có thể thực hiện trong \(O(1)\). Vì vậy phần
sau chỉ xét bài toán gồm \(N\) hình chữ nhật trên lưới có tọa độ nguyên từ
\(1\) đến \(2N\).

\subsection{Một ý tưởng thô}

Cách thô nhất giống thuật toán loang màu. Dùng một mảng \(2N\times 2N\) ghi
màu hiện tại của từng ô rời rạc, ban đầu chưa có màu. Xử lý các hình chữ nhật
từ trên xuống dưới; với mỗi hình, tô màu cho phần bên dưới nó mà vẫn chưa có
màu.

Nếu không đủ bộ nhớ cho mảng hai chiều, có thể xử lý từng hàng rời rạc. Với
mỗi hàng, dùng mảng một chiều độ dài \(2N\). Ta lại quét các hình chữ nhật từ
trên xuống dưới; nếu hình chữ nhật cắt hàng hiện tại, tô các ô chưa màu trong
đoạn giao của nó với hàng đó. Cách này hầu như không tăng thời gian, chỉ tăng
một hằng số nhỏ.

\section{Giải quyết bài toán}

\subsection{Thuật toán kinh điển}

Lời giải kinh điển cho dạng bài này là cây đoạn. Với mỗi hàng rời rạc, ta
dùng một cấu trúc cây \(O(N)\) bộ nhớ. Cây hỗ trợ thao tác gán thuộc tính
\((l,r,X)\), nghĩa là gán thuộc tính \(X\) cho đoạn \([l,r]\).

Thao tác chuẩn trên nút \(P\) có thể mô tả như sau:

\begin{enumerate}
  \item Nếu \([l,r]\) không giao với đoạn của \(P\), trả về.
  \item Nếu \([l,r]\) phủ toàn bộ đoạn của \(P\), ghi thuộc tính \(X\) tại
  nút này.
  \item Ngược lại, đệ quy xuống hai con của \(P\).
\end{enumerate}

Với bài toán hiện tại, thao tác cần hơi đặc biệt: chèn một đoạn màu \(X\)
\([l,r]\), nhưng các phần đã có màu không được thay thế; sau khi xử lý xong
một hàng, cần trả về tổng độ dài đang bị mỗi màu phủ.

Thao tác thống kê màu chỉ cần duyệt toàn bộ cây, mất \(O(N)\). Thao tác chèn
được sửa như sau:

\begin{enumerate}
  \item Nếu \([l,r]\) không giao với đoạn của \(P\), trả về.
  \item Nếu nút \(P\) đã bị một màu nào đó phủ, trả về.
  \item Nếu \([l,r]\) phủ toàn bộ đoạn của \(P\), ghi màu \(X\) tại nút này.
  \item Ngược lại, đệ quy xuống hai con của \(P\).
\end{enumerate}

Thuật toán này dùng \(O(N)\) bộ nhớ và thời gian
\[
  O(N^2\log N),
\]
và có thể xem là một lời giải tốt cho dữ liệu bài toán.

\subsection{Một thuật toán khác}

Bỏ qua lối mòn cây đoạn, ta quay lại điểm xuất phát và phân tích thuật toán
thô ở phần trước. Một bài toán thường có hai con đường: dùng cấu trúc dữ liệu
cao cấp, hoặc xem lại thuật toán ban đầu, tìm phần dư thừa rồi cải tiến đúng
chỗ.

Điểm dư thừa rất rõ: khi xử lý một đoạn hiện tại, nếu gặp vùng đã được phủ,
thuật toán thô không thể nhảy qua nhanh mà vẫn phải so sánh từng ô. Ta lập
mảng \(\mathit{next}\), ban đầu \(\mathit{next}[i]=i+1\), và mảng
\(\mathit{color}\), trong đó \(\mathit{color}[i]\) ghi màu của ô rời rạc thứ
\(i\) trên hàng hiện tại. Khi chèn đoạn \([l,r)\), đoạn này chiếm các ô
\(l,l+1,\ldots,r-1\).

Thuật toán:
\begin{enumerate}
  \item Đặt con trỏ \(i=l\).
  \item Nếu \(\mathit{color}[i]\) chưa có màu, tô nó bằng màu hiện tại.
  \item Di chuyển \(i\) đến \(\mathit{next}[i]\), đồng thời ghi lại các vị
  trí con trỏ đã đi qua.
  \item Nếu \(i\le r-1\), quay lại bước 2.
  \item Sau khi kết thúc, đặt \(\mathit{next}[x]\) của mọi vị trí \(x\) đã đi
  qua thành \(\mathit{next}[r-1]\).
\end{enumerate}

Ví dụ, sau khi phủ màu \(1\) trên \([2,6)\), các ô \(2,3,4,5\) được tô và các
con trỏ tương ứng được nhảy thẳng đến \(6\). Khi phủ màu \(2\) trên \([4,7)\),
con trỏ bắt đầu tại \(4\), nhận ra vùng đã tô, nhảy qua đến \(6\), chỉ tô ô
\(6\), rồi dừng tại \(7\). Khi phủ màu \(3\) trên \([1,9)\), con trỏ chỉ thật
sự tô các ô còn trống và nhảy qua các cụm đã có màu.

Tinh túy của thuật toán là hợp nhất các ô rời rạc đã được tô liền kề, hoặc
đánh dấu trước để lần sau chúng tự động được hợp nhất. Đây chính là tư tưởng
phân tích khấu hao, gợi nhớ đến rừng tập rời nhau có nén đường đi và hợp nhất
theo hạng. Tuy nhiên phiên bản trên chưa nén đường đi đầy đủ, cũng chưa hợp
nhất theo hạng, nên còn có thể hoàn thiện.

\subsection{Hoàn thiện bằng nén đường đi đầy đủ}

Xem các đoạn đã tô kề nhau là một tập hợp. Khi tìm kiếm trong \([l,r)\), nếu
gặp đoạn đã tô, ta muốn nhảy qua cả tập hợp đó trong thời gian khấu hao
\(O(1)\).

Bổ sung mảng \(p\), ban đầu \(p[i]=i\), giống con trỏ cha trong cấu trúc
disjoint-set. Mảng \(\mathit{next}\) vẫn khởi tạo bằng \(\mathit{next}[i]=i+1\),
và mảng \(\mathit{color}\) vẫn ghi màu của ô thứ \(i\).

Ta định nghĩa:
\begin{itemize}
  \item \(i\) là nút đại diện cho ô rời rạc thứ \(i\);
  \item \(p[i]\) là cha của nút \(i\);
  \item nút có \(p[i]=i\) là gốc của một cây.
\end{itemize}

Với nút không phải gốc, giá trị \(\mathit{next}[i]\) không còn ý nghĩa. Với
gốc, ở phiên bản chưa xét hạng, ta có \(\mathit{next}[i]=i+1\).

Thuật toán chèn đoạn được viết lại:
\begin{enumerate}
  \item Đặt \(i=l\).
  \item Nếu \(\mathit{color}[i]\) chưa có màu, tô màu cho nó.
  \item Ghi lại vị trí hiện tại của con trỏ.
  \item Nếu \(i>r-1\), chuyển sang bước cuối.
  \item Nếu \(p[i]\ne i\), đặt \(i=p[i]\); ngược lại đặt
  \(i=\mathit{next}[i]\).
  \item Quay lại bước 2.
  \item Nén đường đi: với mọi vị trí \(x\) con trỏ đã đi qua, đặt
  \(p[x]\) trỏ về đại diện mới của tập hợp sau khi hợp nhất.
\end{enumerate}

Ví dụ cùng các đoạn \([2,6)\), \([4,7)\), \([1,9)\): sau lần phủ đầu, các nút
\(2\) đến \(5\) được hợp nhất, gốc là \(5\). Lần phủ thứ hai bắt đầu tại \(4\),
theo cha đến \(5\), rồi nhảy đến \(6\), tô \(6\), sau đó hợp nhất hai tập với
gốc mới là \(6\). Lần phủ thứ ba tô \(1\), nhảy qua tập bắt đầu ở \(2\), rồi
tô \(7,8\) và hợp nhất toàn bộ các tập liên quan.

\subsection{Thiết lập hạng}

Hạng ở đây có thể hiểu là số nút trong một tập. Khi hợp nhất hai tập, ta gắn
cây nhỏ hơn vào cây lớn hơn; cùng với nén đường đi, độ phức tạp khấu hao của
các thao tác tập rời nhau là gần \(O(1)\). Trong thuật toán của bài này, chỉ
cần sửa phần chọn gốc mới sau khi hợp nhất: chọn gốc của cây có hạng lớn nhất
làm đại diện mới.

Khi dùng hạng, không còn nhất thiết \(\mathit{next}[i]=i+1\) đối với gốc, nên
cần duy trì riêng giá trị ``vị trí tiếp theo sau cả tập''. Có nhiều nhất
\(2N\) tập, số lần hợp nhất không vượt quá \(2N-1\). Vì vậy thời gian xử lý
một hàng rời rạc là \(O(N)\), và toàn bộ thuật toán đạt
\[
  O(N^2).
\]

\subsection{Tiểu kết}

Từ bỏ cây đoạn, quay về thuật toán thô rồi cải tiến nó là điều không dễ. Quá
trình hình thành thuật toán cũng rất chặt chẽ: trước hết tìm chỗ dư thừa để
đạt cải thiện về lượng; sau khi thấy ưu điểm của cách làm, liên hệ nó với một
cấu trúc dữ liệu tương tự, hấp thụ những ưu điểm như nén đường đi và hợp nhất
theo hạng; cuối cùng đạt được thay đổi về chất trong độ phức tạp.

\section{Kết luận}

Trong truyền thuyết Hy Lạp, vua Phrygia thắt nút Gordius và tiên đoán rằng ai
tháo được nút ấy sẽ thống trị châu Á. Alexander Đại đế không lần theo từng
sợi dây, mà rút kiếm chém đứt nút. Có người nói đó là ``gian lận'', nhưng có
lẽ điều các vị thần cần tìm chính là sự quyết đoán ít lời, nhiều hành động
như vậy.

Đối mặt với một thế giới đa dạng, ta thường rơi vào trạng thái ``ý thức về
đáp án quá mạnh'': nhanh chóng tưởng tượng ra một đáp án, vui vẻ đi theo con
đường đó, nhưng không chú ý xem giả thiết ban đầu có đúng và đầy đủ hay không.
Bị chính chướng ngại vô hình do mình tạo ra che mắt là điều đáng tiếc. Vì vậy
ta không nên dừng bước vì nhận thức phiến diện hoặc vì sợ bỏ lợi ích trước
mắt; đôi khi cần sự quả cảm để quay lại điểm xuất phát và định nghĩa lại vấn
đề. Đó chính là sự thống nhất giữa tính tiến lên và tính quanh co của tư duy.

Cách biểu diễn vấn đề thường quan trọng hơn đáp án. Đáp án có thể chỉ là toán
học hoặc thực nghiệm, còn việc đặt ra câu hỏi mới, khả năng mới, hoặc nhìn
một vấn đề cũ từ góc độ mới mới đòi hỏi trí tưởng tượng sáng tạo. Trở về điểm
xuất phát để định nghĩa lại vấn đề chính là một nguồn gốc của tiến bộ khoa
học.

\section*{Tài liệu tham khảo}

\begin{itemize}
  \item Thomas H. Cormen, Charles E. Leiserson, Ronald L. Rivest, Clifford
  Stein, \textit{Introduction to Algorithms}.
  \item Donald E. Knuth, \textit{The Art of Computer Programming}.
  \item Allen B. Tucker et al., \textit{The Computer Science and Engineering
  Handbook}.
  \item Wu Wenhu, Wang Jiande, \textit{Phân tích và thiết kế chương trình cho
  thuật toán thực dụng}.
  \item Lời giải Baltic OI 2004 do Xin Tao cung cấp tại \texttt{http://get-me.to/oi}.
\end{itemize}

\section*{Lời cảm ơn}

Tác giả cảm ơn Xin Tao ở tỉnh Liêu Ninh vì sự giúp đỡ lớn, gồm một phần thuật
toán và chứng minh; cảm ơn các thầy cô trong nhóm huấn luyện Olympic tin học
thanh thiếu niên tỉnh Giang Tô; cảm ơn các thầy cô của Trường Ngoại ngữ Nam
Kinh, gia đình, và tất cả bạn bè trong giới tin học.

\end{document}
